\documentclass{labreport}

\setlength{\headheight}{28pt}
\addtolength{\topmargin}{-13pt}

% UTF-8 support for Romanian diacritics in code listings
\lstset{
  literate={ă}{{\v{a}}}1
           {â}{{\^{a}}}1
           {î}{{\^{i}}}1
           {ș}{{\c{s}}}1
           {ț}{{\c{t}}}1
           {Ă}{{\v{A}}}1
           {Â}{{\^{A}}}1
           {Î}{{\^{I}}}1
           {Ș}{{\c{S}}}1
           {Ț}{{\c{T}}}1
}

% Go language definition for listings
\lstdefinelanguage{Golang}{
  morekeywords=[1]{
    package,import,func,return,type,struct,var,const,interface,
    if,else,switch,case,default,for,range,break,continue,go,select,defer,chan
  },
  morekeywords=[2]{
    string,int,int64,bool,byte,error,map,make,len,cap,append,delete,close,nil
  },
  sensitive=true,
  morecomment=[l]{//},
  morecomment=[s]{/*}{*/},
  morestring=[b]",
  morestring=[b]`
}

\lstdefinestyle{golang}{
  language=Golang,
  basicstyle=\ttfamily\small,
  keywordstyle=[1]\color{blue}\bfseries,
  keywordstyle=[2]\color{teal}\bfseries,
  commentstyle=\color{gray}\itshape,
  stringstyle=\color{red!70!black},
  showstringspaces=false,
  breaklines=true,
  frame=single,
  numbers=left,
  numberstyle=\tiny\color{gray},
  backgroundcolor=\color{gray!10}
}

\begin{document}

% ---------------------------------------------------------
% 1. TITLE PAGE
% ---------------------------------------------------------
\labtitle{2}{Laboratory Work 2: A Coffee Shop Server}{Poiata Calin\\Covali Ilie}{gr. FAF-243}{asist. univ. Cristofor Fistic}

\tableofcontents
\newpage

% ---------------------------------------------------------
% 2. OBJECTIVE AND SYSTEM SPECIFICATION
% ---------------------------------------------------------
\section{Objective and Architectural Requirements}

\subsection{Laboratory Objective}
The objective of this laboratory work is to design and build a concurrent, fault-tolerant network server implementing a simulated coffee shop using raw TCP sockets directly. Clients connect across the network to view stock and order products, which are prepared concurrently by a pool of workers matching hardware CPU capacity. A provider connects to supply inventory and receive low-stock alerts. 

Key architectural requirements include:
\begin{enumerate}
    \item \textbf{Direct Socket API Usage:} Utilizing raw OS-level TCP socket primitives (\texttt{socket}, \texttt{bind}, \texttt{listen}, \texttt{accept}, \texttt{connect}, \texttt{send}/\texttt{write}, \texttt{recv}/\texttt{read}, and \texttt{close}). High-level wrappers such as HTTP/REST, WebSockets, RPC, or messaging frameworks are strictly prohibited.
    \item \textbf{Stream Framing and Message Boundary Reconstruction:} Implementing explicit buffer accumulation and delimiter slicing to handle TCP stream fragmentation, split packets, and batched deliveries.
    \item \textbf{Worker Thread Pool and Shared State Protection:} Creating a fixed-size worker pool pinned to the number of logical CPU cores (\texttt{runtime.NumCPU()}), consuming orders through a FIFO queue, with shared inventory protected against race conditions and overselling.
    \item \textbf{Asynchronous Push Notifications:} Communicating order readiness and stock threshold alerts asynchronously without client-side polling.
    \item \textbf{Fault Tolerance and Dynamic State Rollback:} Detecting sudden client failure (crash, SIGKILL, network drop), cancelling queued or in-flight preparations, and atomically returning reserved goods to inventory.
    \item \textbf{Heterogeneous Multi-Language Bonus (Section 8):} Developing the server engine in \textbf{Go} (Member 1) and the interactive client, provider, and stress-test suite in \textbf{Python} (Member 2), operating seamlessly over raw TCP.
\end{enumerate}

\subsection{Hardware and Execution Environment}
The server and test harness were deployed and benchmarked on Arch Linux. Hardware and environment parameters are summarized in Table~\ref{tab:hardware}.

\begin{table}[H]
\centering
\caption{System and Hardware Environment}
\label{tab:hardware}
\begin{tabular}{ll}
\toprule
\textbf{Parameter} & \textbf{Specification} \\
\midrule
CPU Model & AMD Ryzen 7 4700U with Radeon Graphics \\
CPU Architecture & x86\_64 (Zen 2, 8 Physical Cores, 8 Logical Processors) \\
RAM & 16 GB DDR4 \\
Operating System & Arch Linux (Linux Kernel 7.2.11-arch1-1 / Rolling) \\
Server Implementation & Go 1.27 (Native goroutines, channels, \texttt{sync.Mutex}) \\
Client \& Test Harness & Python 3.14 (Raw \texttt{socket}, dual-threading, \texttt{unittest}) \\
Network Transport & Raw TCP Socket (Configurable port, default 8080) \\
Framing Format & Newline-delimited text (\texttt{\textbackslash n}), UTF-8 encoded \\
\bottomrule
\end{tabular}
\end{table}

\subsection{Workload Breakdown and Role Distribution}
The project was executed through an exact 50/50 division of responsibility between the two team members, as outlined in Table~\ref{tab:distribution}.

\begin{table}[H]
\centering
\caption{Team Responsibility and Task Matrix}
\label{tab:distribution}
\resizebox{\textwidth}{!}{%
\begin{tabular}{lll}
\toprule
\textbf{Component} & \textbf{Member 1 (Go)} & \textbf{Member 2 (Python)} \\
\midrule
Primary Artifacts & \texttt{go/server/}, \texttt{go/main.go} & \texttt{python/client.py}, \texttt{provider.py}, \texttt{test\_harness/} \\
Core Architecture & TCP Listener, Worker Pool, Mutexes & REPL CLIs, Dual-Threaded I/O, Chaos Harness \\
Tasks Covered & Server-side Tasks 1, 2, 3, 4 & Client-side Tasks 1, 3, 4 + Automated Defense Suite \\
Key Challenge & Channel queuing, atomic rollback, context & Non-blocking stdin, barrier sync, protocol fuzzer \\
\bottomrule
\end{tabular}%
}
\end{table}

\newpage

% ---------------------------------------------------------
% 3. WIRE PROTOCOL SPECIFICATION
% ---------------------------------------------------------
\section{Wire Protocol Specification}

\subsection{Framing and Transport}
Because TCP is an unformatted continuous stream of bytes without record boundaries, our protocol enforces explicit line framing:
\begin{equation}
\text{Message} = \text{Payload} \mathbin{\Vert} \texttt{"\textbackslash n"}
\end{equation}
Every request and response terminates with a single newline character (\texttt{\textbackslash n} or \texttt{\textbackslash r\textbackslash n}). The encoding is UTF-8.

\subsection{Connection Handshake}
Upon establishing a TCP connection, the remote party must identify its role before issuing any operational commands:
\begin{itemize}
    \item \textbf{Client identification:} Sends \texttt{ROLE client\textbackslash n}. Server replies: \texttt{OK role\_accepted\textbackslash n}.
    \item \textbf{Provider identification:} Sends \texttt{ROLE provider\textbackslash n}. Server replies: \texttt{OK role\_accepted\textbackslash n}.
    \item \textbf{Invalid handshake:} Any other message results in \texttt{ERROR expected\_role\_handshake\textbackslash n} or \texttt{ERROR invalid\_role\textbackslash n}, followed by immediate socket teardown.
\end{itemize}

\subsection{Command and Notification Grammar}
Table~\ref{tab:protocol} details the full wire protocol command set.

\begin{table}[H]
\centering
\caption{TCP Wire Protocol Command Grammar}
\label{tab:protocol}
\resizebox{\textwidth}{!}{%
\begin{tabular}{lll}
\toprule
\textbf{Command / Event} & \textbf{Direction} & \textbf{Format / Example} \\
\midrule
\texttt{ROLE} & Client/Provider $\to$ Srv & \texttt{ROLE client} / \texttt{ROLE provider} \\
\texttt{OK} & Server $\to$ Peer & \texttt{OK role\_accepted} \\
\texttt{GET\_MENU} & Client/Provider $\to$ Srv & \texttt{GET\_MENU} \\
\texttt{MENU} & Server $\to$ Peer & \texttt{MENU coffee=35 tea=20 croissant=15} \\
\texttt{ORDER} & Client $\to$ Server & \texttt{ORDER coffee=2 tea=1} \\
\texttt{ORDER\_ACCEPTED} & Server $\to$ Client & \texttt{ORDER\_ACCEPTED id=1 wait\_time=11} \\
\texttt{ORDER\_REFUSED} & Server $\to$ Client & \texttt{ORDER\_REFUSED reason=insufficient\_stock coffee\_left=1 ...} \\
\texttt{ORDER\_READY} & Server $\to$ Client & \texttt{ORDER\_READY id=1} \textit{(Asynchronous Push)} \\
\texttt{DELIVER} & Provider $\to$ Server & \texttt{DELIVER coffee=20 croissant=10} \\
\texttt{DELIVERY\_OK} & Server $\to$ Provider & \texttt{DELIVERY\_OK coffee=55 tea=20 croissant=25} \\
\texttt{ALERT\_LOW\_STOCK} & Server $\to$ Provider & \texttt{ALERT\_LOW\_STOCK product=croissant count=2} \textit{(Push)} \\
\texttt{ERROR} & Server $\to$ Peer & \texttt{ERROR <error\_code\_or\_description>} \\
\bottomrule
\end{tabular}%
}
\end{table}

\newpage

% ---------------------------------------------------------
% 4. TASK 1: SINGLE-CLIENT CORE AND STREAM FRAMING
% ---------------------------------------------------------
\section{Task 1: Single-Client Architecture and Stream Framing}

\subsection{The TCP Byte-Stream Problem}
A fundamental concept in network systems is that TCP provides a reliable, ordered \textit{byte stream}, not a message stream. Network layers, OS TCP socket buffers, and network interfaces may segment or coalesce writes arbitrarily:
\begin{enumerate}
    \item \textbf{Packet Splitting (Partial Message):} A client issues a single write of \texttt{"ORDER coffee=1\textbackslash n"}, but the network card transmits the first 8 bytes in packet 1 and the remainder in packet 2. The server must not fail with an unparseable token.
    \item \textbf{Packet Coalescing (Multiple Messages in One Read):} A client rapidly sends \texttt{"ROLE client\textbackslash n"} and \texttt{"GET\_MENU\textbackslash n"}. The server's OS socket buffer coalesces these into a single 4096-byte chunk. If the application reads without framing, both commands are incorrectly merged.
\end{enumerate}

\subsection{Framing Buffer Implementation}
To resolve this, the Go server implements a buffered reader in \texttt{Connection.ReadLine()} (\texttt{go/server/connection.go}). An internal accumulator (\texttt{strings.Builder}) stores unparsed bytes between reads.

\begin{lstlisting}[style=golang, caption={Newline Delimiter Slicing in Connection.ReadLine()}]
func (c *Connection) ReadLine() (string, error) {
    c.readMu.Lock()
    defer c.readMu.Unlock()

    if c.isClosed {
        return "", io.EOF
    }
    buf := make([]byte, 4096)

    for {
        current := c.readBuf.String()
        // 1. Check if complete newline exists in buffer
        if idx := strings.IndexByte(current, '\n'); idx >= 0 {
            line := current[:idx]
            remainder := current[idx+1:]
            c.readBuf.Reset()
            c.readBuf.WriteString(remainder)
            return strings.TrimRight(line, "\r"), nil
        }

        // 2. Buffer does not contain full message; read chunk from socket
        n, err := c.conn.Read(buf)
        if n > 0 {
            c.readBuf.Write(buf[:n])
        }
        if err != nil {
            // EOF or client disconnection
            remaining := c.readBuf.String()
            c.readBuf.Reset()
            if len(strings.TrimSpace(remaining)) > 0 {
                return strings.TrimRight(remaining, "\r"), nil
            }
            return "", err
        }
    }
}
\end{lstlisting}

\subsection{Immediate Acknowledgment vs. Asynchronous Delivery}
When an order arrives, the preparation time is calculated as the sum of its item preparation durations:
\begin{equation}
T_{\text{wait}} = \sum_{p \in \text{Items}} (\text{PrepTime}(p) \times \text{Quantity}(p))
\end{equation}
where $\text{PrepTime}(\text{coffee}) = 4\text{s}$, $\text{PrepTime}(\text{tea}) = 3\text{s}$, and $\text{PrepTime}(\text{croissant}) = 2\text{s}$. 

Rather than blocking the client connection synchronously for $T_{\text{wait}}$ seconds, the server immediately deducts the inventory and dispatches \texttt{ORDER\_ACCEPTED id=<id> wait\_time=<sec>}. This non-blocking architecture allows the client to keep reading the menu or placing additional orders while previous orders are prepared in the background.

\newpage

% ---------------------------------------------------------
% 5. TASK 2: MULTI-CLIENT SERVING AND CPU WORKER POOL
% ---------------------------------------------------------
\section{Task 2: Multi-Client Serving and CPU Worker Pool}

\subsection{Concurrency Model: Goroutines vs. OS Threads}
Traditional server architectures allocate one OS thread per client connection. With 1,000 concurrent clients, allocating $1\text{ MB}$ to $8\text{ MB}$ stack memory per thread would consume gigabytes of RAM and trigger severe OS scheduler contention.

In our Go server, client handling uses \textbf{Goroutines}. Go's M:N runtime multiplexes thousands of goroutines across 8 OS threads (matching physical CPU cores). Each goroutine starts with only $2\text{ KB}$ of contiguous stack space, which grows dynamically on demand. The listener loop spawns a goroutine per connection:
\begin{lstlisting}[style=golang]
for {
    conn, err := listener.Accept()
    if err != nil { continue }
    c := NewConnection(conn)
    go s.handleConnection(c)
}
\end{lstlisting}

\subsection{Thread-Safe Inventory Management}
The coffee shop stock is shared across all concurrent connections and workers. If two clients order the last available croissant at the exact same instant, an unprotected read-modify-write sequence would cause both orders to be accepted, driving the inventory to $-1$.

To prevent this, \texttt{Stock} encapsulates an internal \texttt{sync.Mutex} (\texttt{go/server/stock.go}). The method \texttt{TryDeduct()} executes inventory verification and deduction as a single atomic critical section:

\begin{lstlisting}[style=golang, caption={Atomic Stock Verification and Deduction}]
func (s *Stock) TryDeduct(requested map[string]int) (bool, map[string]int) {
    s.mu.Lock()
    defer s.mu.Unlock()

    // 1. Verify all requested items are in stock
    for item, qty := range requested {
        available, exists := s.items[item]
        if !exists || available < qty {
            return false, s.snapshotLocked()
        }
    }

    // 2. Atomically deduct requested quantities
    for item, qty := range requested {
        s.items[item] -= qty
    }
    return true, s.snapshotLocked()
}
\end{lstlisting}

\subsection{CPU Worker Pool and FIFO Channel Queue}
In compliance with Section 2 of the specification, the server creates exactly one worker for each logical processor:
\begin{equation}
N_{\text{workers}} = \texttt{runtime.NumCPU()} = 8
\end{equation}

Workers consume orders from a Go channel (\texttt{s.orderQueue chan *Order}), which acts as a thread-safe, lock-free FIFO queue. When an order is accepted, it is enqueued into the channel. Free workers pull the oldest waiting order and simulate preparation.

\begin{lstlisting}[style=golang, caption={Worker Loop with Cancellation and Readiness Dispatch}]
func (s *Server) workerLoop(workerID int) {
    for order := range s.orderQueue {
        if order.Ctx.Err() != nil {
            s.removeClientOrder(order)
            continue
        }

        select {
        case <-time.After(order.TotalTime):
            order.Mu.Lock()
            order.Completed = true
            order.Mu.Unlock()
            // Asynchronously dispatch notification to client
            _ = order.Conn.WriteLine(fmt.Sprintf("ORDER_READY id=%d", order.ID))
        case <-order.Ctx.Done():
            // Aborted mid-prep due to client disconnect
        }
        s.removeClientOrder(order)
    }
}
\end{lstlisting}

\subsection{Full-Duplex Socket Concurrency Protection}
While a worker is preparing an order, the client's goroutine may concurrently execute \texttt{GET\_MENU} and write to the same TCP socket. Because simultaneous writes from two goroutines could interleave bytes and corrupt messages, \texttt{Connection.WriteLine()} serializes writes via an internal \texttt{writeMu sync.Mutex}.

\newpage

% ---------------------------------------------------------
% 6. TASK 3: PROVIDER RESTOCKING AND ASYNCHRONOUS ALERTS
% ---------------------------------------------------------
\section{Task 3: Provider Restocking and Low-Stock Alerts}

\subsection{Role Discrimination and Restocking}
When a connection identifies as \texttt{ROLE provider}, the server stores a reference in \texttt{s.providerConn}. If a client attempts to invoke \texttt{DELIVER}, the server returns \texttt{ERROR unauthorized\_role}. Conversely, deliveries from a provider invoke \texttt{s.stock.AddStock(items)}, atomically replenishing inventory:
\begin{equation}
\text{Stock}_{\text{new}}(p) = \text{Stock}_{\text{old}}(p) + \text{Delivered}(p), \quad \forall p \in \{\text{coffee}, \text{tea}, \text{croissant}\}
\end{equation}

\subsection{Threshold-Triggered Low-Stock Alerts}
Section 5 specifies that whenever any product falls below 5 units, the provider must be alerted immediately.
\begin{itemize}
    \item \textbf{Online Provider:} The server dispatches \texttt{ALERT\_LOW\_STOCK product=<p> count=<qty>\textbackslash n} asynchronously over the provider's active TCP socket.
    \item \textbf{Offline Provider:} The alert is enqueued in \texttt{s.pendingAlerts}. When a provider eventually connects and passes the handshake, the server flushes all queued alerts before entering the command loop.
\end{itemize}

\begin{figure}[H]
\centering
\begin{verbatim}
       Client Orders (Stock drops < 5)
                     |
                     v
          [checkLowStockAndAlert()]
                     |
           +---------+---------+
           |                   |
    Provider Online?    Provider Offline?
           |                   |
           v                   v
   Push ALERT_LOW_STOCK    Queue in s.pendingAlerts
   over TCP socket         Flush upon ROLE provider connect
\end{verbatim}
\caption{Dual-Mode Stock Alert Dispatch Workflow}
\label{fig:alerts}
\end{figure}

\newpage

% ---------------------------------------------------------
% 7. TASK 4: FAULT TOLERANCE AND CLEANUP
% ---------------------------------------------------------
\section{Task 4: Fault Tolerance, Abrupt Disconnection, and State Rollback}

\subsection{Disconnection Detection}
A client process can terminate abruptly due to a kill signal (\texttt{SIGKILL}), terminal closure, or network loss.
\begin{itemize}
    \item \textbf{Clean Termination / OS Kill:} The client operating system sends a TCP \texttt{FIN} or \texttt{RST} packet. In Go, \texttt{c.conn.Read()} immediately returns $0$ bytes (\texttt{io.EOF}) or a connection reset error. The server detects this in \texttt{ReadLine()} and initiates cleanup.
    \item \textbf{Silent Network Drop:} Physical link severance does not send packets. The connection remains open until a socket write fails with a broken pipe (\texttt{EPIPE}) or TCP keep-alive probes time out.
\end{itemize}

\subsection{Order Cancellation and Stock Rollback Invariant}
When a client disconnects unexpectedly, the server must satisfy strict state conservation:
\begin{equation}
\text{Stock}_{\text{after}} = \text{Stock}_{\text{before}} + \sum_{o \in \text{UncompletedOrders}} o.\text{Items}
\end{equation}

To implement this:
\begin{enumerate}
    \item Every order is created with a cancellation context:
    \begin{lstlisting}[style=golang]
    ctx, cancel := context.WithCancel(context.Background())
    \end{lstlisting}
    \item Active orders are tracked in a thread-safe registry: \texttt{s.activeOrders[clientID]}.
    \item Upon client disconnection, \texttt{s.cleanupClient(clientID)} is called. It invokes \texttt{order.Cancel()} for all uncompleted orders.
    \item If the order is queued, the worker discards it upon dequeuing. If the order is currently being prepared, the worker's \texttt{select} statement immediately wakes up on \texttt{<-order.Ctx.Done()}, aborting preparation.
    \item All reserved goods from uncompleted orders are returned to inventory via \texttt{s.stock.RestoreStock(restoredStock)}.
\end{enumerate}

\begin{lstlisting}[style=golang, caption={Atomic Cleanup and Stock Rollback on Disconnect}]
func (s *Server) cleanupClient(clientID string) {
    s.ordersMu.Lock()
    clientOrders, exists := s.activeOrders[clientID]
    if exists { delete(s.activeOrders, clientID) }
    s.ordersMu.Unlock()

    if !exists || len(clientOrders) == 0 { return }

    restoredStock := make(map[string]int)
    for order := range clientOrders {
        order.Mu.Lock()
        if !order.Completed {
            order.Cancel() // Interrupt worker goroutine
            for prod, qty := range order.Items {
                restoredStock[prod] += qty
            }
        }
        order.Mu.Unlock()
    }

    if len(restoredStock) > 0 {
        s.stock.RestoreStock(restoredStock)
        log.Printf("[%s] Disconnected. Restored stock: %v\n", clientID, restoredStock)
    }
}
\end{lstlisting}

\newpage

% ---------------------------------------------------------
% 8. EMPIRICAL VERIFICATION AND STRESS TESTING
% ---------------------------------------------------------
\section{Empirical Verification and Test Suite Results}

To validate all requirements, Member 2 constructed an automated Python test harness targeting the Go server.

\subsection{High-Concurrency Race Condition Benchmark}
The script \texttt{test\_race.py} evaluates synchronization under extreme contention. 
\begin{itemize}
    \item \textbf{Setup:} Stock of croissants was queried ($S = 17$). A barrier synchronized $N = 27$ concurrent client threads ($N = S + 10$).
    \item \textbf{Execution:} All 27 threads simultaneously sent \texttt{ORDER croissant=1}.
    \item \textbf{Result:} Exactly 17 orders were accepted, and exactly 10 orders were refused. Final croissant stock was exactly 0. Zero race conditions or overselling occurred.
\end{itemize}

\begin{table}[H]
\centering
\caption{Task 2 Race Condition Test Results}
\label{tab:race}
\begin{tabular}{ll}
\toprule
\textbf{Metric} & \textbf{Value} \\
\midrule
Total Concurrent Client Threads & 27 \\
Initial Croissant Inventory & 17 \\
Orders Accepted & 17 (100\% available) \\
Orders Refused (Insufficient Stock) & 10 \\
Overselling / Stock Discrepancies & 0 \\
Final Croissant Inventory & 0 (Perfect invariant conservation) \\
\textbf{Status} & \textbf{PASS} \\
\bottomrule
\end{tabular}
\end{table}

\subsection{Fault Tolerance and Rollback Test}
The script \texttt{test\_cancellation.py} validates client crash recovery.
\begin{itemize}
    \item \textbf{Setup:} Initial coffee stock was $50$. A client ordered 3 coffees (wait time $12\text{s}$). Stock became $47$.
    \item \textbf{Fault Injection:} At elapsed time $t = 1.5\text{s}$, the client process closed its TCP socket abruptly.
    \item \textbf{Result:} The Go server detected \texttt{io.EOF}, interrupted Worker 4 mid-preparation, discarded the elapsed work, and restored the 3 coffees. Independent query verified final coffee stock returned to exactly $50$.
\end{itemize}

\subsection{Protocol Fuzzing and Resilience Suite}
The test \texttt{fuzzer.py} subjected the Go server to hostile and corrupted TCP payloads:
\begin{itemize}
    \item \textbf{Oversized Buffers:} 64 KB single-line payloads were handled without memory leaks or buffer overflow.
    \item \textbf{Random Binary Data:} High-entropy raw random bytes were rejected with \texttt{ERROR unknown\_command}.
    \item \textbf{Malformed Parameters:} Negative numbers (\texttt{coffee=-4}), zero counts, and unknown products (\texttt{beer}) received strict protocol error responses without hanging other threads.
    \item \textbf{Server Survivability:} Post-fuzzing health checks confirmed the server remained fully operational.
\end{itemize}

\newpage

% ---------------------------------------------------------
% 9. CONTROL QUESTIONS
% ---------------------------------------------------------
\section{Control Questions}

\begin{enumerate}
    \item \textbf{What is the difference between the listening socket and the sockets returned by \texttt{accept}? How many of each does your server have with five clients and the provider?}\\
    The \textbf{listening socket} is bound to a well-known local port (e.g., port 8080) in the \texttt{LISTEN} state. Its sole purpose is to receive incoming TCP connection requests (SYN packets) and complete the three-way handshake; it never exchanges application data. The sockets returned by \textbf{\texttt{accept}} are dedicated, connected sockets in the \texttt{ESTABLISHED} state, uniquely bound to a specific client session and used for bidirectional payload transfer (\texttt{read}/\texttt{write}).\\
    With 5 clients and 1 provider, the server has \textbf{1 listening socket} and \textbf{6 accepted sockets} (Total = 7 open socket file descriptors).

    \item \textbf{Why is one send not one receive? How does your protocol deal with it?}\\
    TCP is a streaming protocol operating over variable-sized IP datagrams. It guarantees byte order and reliability, but does not preserve application message boundaries. Senders may coalesce small writes into a single segment (Nagle's algorithm), or intermediate network routers may fragment a single write across multiple IP packets. Consequently, three separate \texttt{send()} calls may arrive in one \texttt{recv()} call, or one \texttt{send()} call may require two \texttt{recv()} calls to read completely. Our protocol deals with this via \textbf{newline-delimited text framing (\texttt{\textbackslash n})}: the server accumulates bytes in an internal buffer (\texttt{readBuf}) and slices out messages strictly at newline boundaries.

    \item \textbf{What does it mean when receive returns 0? And when it fails because the connection was reset?}\\
    When receive returns 0 bytes (in Go, \texttt{io.EOF}), it signifies an \textbf{orderly connection termination} initiated by the peer (the remote OS transmitted a TCP \texttt{FIN} packet indicating it has finished writing data). When receive fails with a \textbf{connection reset} (\texttt{ECONNRESET} or TCP \texttt{RST} packet), it indicates an \textbf{abrupt, abnormal failure}---such as the peer process crashing, being forcefully killed with \texttt{SIGKILL}, or attempting to send data to a closed port.

    \item \textbf{Where in your server is a producer--consumer pattern? Name the producers, the consumers and the queue between them.}\\
    The producer--consumer pattern is implemented in the order processing pipeline:
    \begin{itemize}
        \item \textbf{Producers:} The client connection handler goroutines running \texttt{handleOrder()}. Upon validating that inventory is available, they produce \texttt{Order} objects.
        \item \textbf{Queue:} The Go buffered channel \texttt{s.orderQueue chan *Order}, providing thread-safe FIFO queuing.
        \item \textbf{Consumers:} The fixed pool of 8 worker goroutines running \texttt{workerLoop()}, pulling orders from the channel and simulating preparation.
    \end{itemize}

    \item \textbf{What does it cost when a thread waits for a client? What would change with 1000 clients?}\\
    In operating systems using traditional 1:1 kernel threads (e.g., in C, Java, or standard Python), each thread allocates a static stack of $1\text{ MB}$ to $8\text{ MB}$ in virtual memory. 1,000 idle waiting threads would consume multiple gigabytes of memory, and the OS kernel scheduler would incur severe CPU cache invalidation and context-switching overhead.\\
    In our Go server, connections are handled by \textbf{goroutines} scheduled in user-space by Go's M:N runtime scheduler. A waiting goroutine starts with only $2\text{ KB}$ of memory and does not hold an OS thread while blocked on socket I/O (epoll). 1,000 waiting clients consume only $\approx 2\text{ MB}$ of RAM with virtually zero CPU scheduler penalty.

    \item \textbf{Explain step by step a race condition that could happen on your server without its locks.}\\
    Suppose the remaining inventory of croissants is $1$. Two clients, Client A and Client B, issue an order for 1 croissant simultaneously on separate CPU cores:
    \begin{enumerate}
        \item \textbf{Read Step:} Core 1 (Client A) reads \texttt{stock["croissant"]} from memory into its register; it reads $1$. Concurrently, Core 2 (Client B) reads \texttt{stock["croissant"]} into its register; it also reads $1$.
        \item \textbf{Evaluation Step:} Both cores evaluate \texttt{available >= 1} as \texttt{true}.
        \item \textbf{Write Step:} Core 1 decrements its register to $0$ and writes $0$ back to memory. Core 2 decrements its register to $0$ and writes $-1$ (or $0$) back to memory.
        \item \textbf{Result:} Both clients receive \texttt{ORDER\_ACCEPTED}, but only 1 physical croissant existed. The shop has oversold inventory, causing data corruption.
    \end{enumerate}

    \item \textbf{What identifies one TCP connection? How can five clients on the same machine all be connected to the same port of the server?}\\
    A TCP connection is uniquely identified by an OS-level \textbf{4-tuple}:
    \begin{equation}
    (\text{Source IP}, \text{Source Port}, \text{Destination IP}, \text{Destination Port})
    \end{equation}
    When five clients on the same local machine connect to the server at \texttt{127.0.0.1:8080}, the Destination IP (\texttt{127.0.0.1}), Destination Port (\texttt{8080}), and Source IP (\texttt{127.0.0.1}) are identical. However, the client OS networking stack assigns each client process a unique, distinct \textbf{ephemeral source port} (e.g., 51234, 51235, 51236, 51237, 51238). Because the source port in the 4-tuple differs for each connection, the kernel de-multiplexes incoming packets to the correct socket buffer without ambiguity.
\end{enumerate}

\newpage

% ---------------------------------------------------------
% 10. CONCLUSION
% ---------------------------------------------------------
\section{Conclusion}

This laboratory work successfully demonstrated the design, concurrent architecture, and fault tolerance of a raw TCP client--server system.

\subsection{Summary of Accomplishments}
\begin{enumerate}
    \item \textbf{Zero-Framework Socket Networking:} Built an operational network server without HTTP, RPC, or third-party wrappers, managing newline stream framing and partial TCP read buffers directly.
    \item \textbf{Go Concurrency Primitives:} Leveraged Go goroutines, channels, and \texttt{sync.Mutex} primitives to construct a lock-protected shared inventory and an 8-worker CPU pool.
    \item \textbf{Asynchronous Notification Architecture:} Designed full-duplex socket handling allowing workers to push readiness notifications and low-stock alerts without blocking command execution.
    \item \textbf{Robust Cancellation and Dynamic Rollback:} Implemented instant worker abortion using Go \texttt{context.Context} and guaranteed atomic stock restoration upon unexpected peer disconnects.
    \item \textbf{Multi-Language Collaboration (Section 8 Bonus):} Integrated the Go server engine with the Python client, provider, and automated test harness, achieving complete interoperability and 100\% pass rates across race-condition, cancellation, and fuzzing benchmarks.
\end{enumerate}

\end{document}
